\par\Needspace{6\baselineskip}\medskip
\hypertarget{native-132-article097}{}\pdfbookmark[2]{97.}{article-097}
\begin{center}97.\end{center}
Ad illustrationem methodi inde ab art. 88 expositae duo exempla art. 87 resumemus: anguli auxiliaris $\omega$ significationem hactenus observatam non esse confundendam cum ea, in qua in artt. 86, 87 acceptum erat idem signum, vix opus erit monuisse.

I. In exemplo primo habemus $f=3^0\,47'\,26''{,}865$, porroque
\[
\begin{gathered}
\log\dfrac{r'}r=9{,}991\,4599,\qquad
\log\operatorname{tang}(45^0+\omega)=9{,}997\,8649\,75,\\
\omega=-8'\,27''{,}006.
\end{gathered}
\]
\noindent Hinc per art. 89

\begin{center}
\begin{tabular}{@{}ll@{}}
\makebox[17em][l]{$\log\sin\tfrac12 f^2$\hfill .\,.\,.\,.\,.\,.\,.\,.\,.\,.} & $7{,}038\,9972$\\[.3em]
\makebox[17em][l]{$\log\cos f$\hfill .\,.\,.\,.\,.\,.\,.\,.\,.\,.\,.\,.} & $9{,}999\,0488$\\[.3em]
\hline
 & $7{,}039\,9484$\\[.3em]
\multicolumn{2}{l}{$=\log0{,}001\,0963\,480$}\\[.3em]
\end{tabular}
\end{center}

\begin{center}
\begin{tabular}{@{}ll@{}}
\makebox[17em][l]{$\log\operatorname{tang}2\omega^2$\hfill .\,.\,.\,.\,.\,.\,.\,.} & $5{,}383\,2428$\\[.3em]
\makebox[17em][l]{$\log\cos f$\hfill .\,.\,.\,.\,.\,.\,.\,.\,.\,.\,.} & $9{,}999\,0488$\\[.3em]
\hline
 & $5{,}384\,1940$\\[.3em]
\multicolumn{2}{l}{$=\log0{,}000\,0242\,211$}\\[.3em]
\end{tabular}
\end{center}

\noindent\SourceLeaf{133}adeoque $l=0{,}001\,1205\,691$, $\tfrac56+l=0{,}834\,4539$. Porro fit $\log kt=9{,}576\,6974$

\begin{center}
\begin{tabular}{@{}ll@{}}
\makebox[17em][l]{$2\log kt$\hfill .\,.\,.\,.\,.\,.\,.\,.\,.\,.\,.\,.\,.\,.} & $9{,}153\,3948$\\[.3em]
\makebox[17em][l]{$\mathrm{C.}\tfrac32\log rr'$\hfill .\,.\,.\,.\,.\,.\,.\,.\,.\,.\,.} & $9{,}020\,5181$\\[.3em]
\makebox[17em][l]{$\mathrm{C.}\log8\cos f^3$\hfill .\,.\,.\,.\,.\,.\,.\,.\,.} & $9{,}099\,7636$\\[.3em]
\hline
\makebox[17em][l]{$\log mm$\hfill .\,.\,.\,.\,.\,.\,.\,.\,.\,.\,.\,.\,.\,.} & $7{,}273\,6765$\\[.3em]
\makebox[17em][l]{$\log(\tfrac56+l)$\hfill .\,.\,.\,.\,.\,.\,.\,.\,.\,.\,.\,.} & $9{,}921\,4023$\\[.3em]
\hline
 & $7{,}352\,2742$\\[.3em]
\end{tabular}
\end{center}

\noindent Est itaque valor approximatus ipsius $h=0{,}002\,2504\,7$, cui in tabula nostra II respondet $\log yy=0{,}002\,1633$. Habetur itaque $\log\dfrac{mm}{yy}=7{,}271\,5132$, sive $\dfrac{mm}{yy}=0{,}001\,8685\,87$, unde per formulam 16 fit $x=0{,}000\,7480\,179$; quamobrem quum $\xi$ per tabulam III omnino insensibilis sit, valores inventi pro $h$, $y$, $x$ correctione non indigent. Iam determinatio elementorum ita se habet:

\begin{center}
\begin{tabular}{@{}ll@{}}
\makebox[17em][l]{$\log x$\hfill .\,.\,.\,.\,.\,.\,.\,.\,.\,.\,.\,.\,.\,.\,.\,.} & $6{,}873\,9120$\\[.3em]
\makebox[17em][l]{$\log\sin\tfrac12 g$\hfill .\,.\,.\,.\,.\,.\,.\,.\,.\,.\,.\,.} & $8{,}436\,9560$\\[.3em]
\end{tabular}
\end{center}

\[
\begin{gathered}
\tfrac12 g=1^0\,34'\,2''{,}0286,\qquad
\tfrac12(f+g)=3^0\,27'\,45''{,}4611,\\
\tfrac12(f-g)=19'\,41''{,}4039.
\end{gathered}
\]
\noindent Quare ad normam formularum 27, 28, 29, 30 habetur

\begin{center}
\begin{tabular}{@{}ll@{}}
\makebox[17em][l]{$\log\operatorname{tang}2\omega$\hfill .\,.\,.\,.\,.\,.\,.\,.\,.\,.\,.} & $7{,}691\,6214_n$\\[.3em]
\makebox[17em][l]{$\log\cos\tfrac12(f+g)$\hfill .\,.\,.\,.\,.\,.\,.\,.} & $9{,}999\,2065$\\[.3em]
\makebox[17em][l]{$\log\cos\tfrac12(f-g)$\hfill .\,.\,.\,.\,.\,.\,.\,.} & $9{,}999\,9929$\\[.3em]
\hline
\makebox[17em][l]{$\log P\sin\tfrac12(F-G)$\hfill .\,.\,.\,.\,.} & $7{,}690\,8279_n$\\[.3em]
\makebox[17em][l]{$\log P\cos\tfrac12(F-G)$\hfill .\,.\,.\,.\,.} & $8{,}781\,0240$\\[.3em]
\hline
$\tfrac12(F-G)=$ & $-4^0\,38'\,41''{,}54$\\[.3em]
$\tfrac12(F+G)=$ & $319\quad21\quad38{,}05$\\[.3em]
\hline
$F=$ & $314\quad42\quad56{,}51$\\[.3em]
$v=$ & $310\quad55\quad29{,}64$\\[.3em]
$v'=$ & $318\quad30\quad23{,}37$\\[.3em]
$G=$ & $324\quad0\quad19{,}59$\\[.3em]
$E=$ & $320\quad52\quad15{,}53$\\[.3em]
$E'=$ & $327\quad8\quad23{,}65$\\[.3em]
\end{tabular}
\end{center}

\begin{center}
\begin{tabular}{@{}ll@{}}
\makebox[17em][l]{$\mathrm{C.}\log\cos2\omega$\hfill .\,.\,.\,.\,.\,.\,.\,.\,.\,.} & $0{,}000\,0052$\\[.3em]
\makebox[17em][l]{$\log\sin\tfrac12(f+g)$\hfill .\,.\,.\,.\,.\,.\,.\,.} & $8{,}781\,0188$\\[.3em]
\makebox[17em][l]{$\log\sin\tfrac12(f-g)$\hfill .\,.\,.\,.\,.\,.\,.\,.} & $7{,}757\,9709$\\[.3em]
\hline
\makebox[17em][l]{$\log Q\sin\tfrac12(F+G)$\hfill .\,.\,.\,.\,.} & $7{,}691\,6143_n$\\[.3em]
\makebox[17em][l]{$\log Q\cos\tfrac12(F+G)$\hfill .\,.\,.\,.\,.} & $7{,}757\,9761$\\[.3em]
\hline
\makebox[17em][l]{$\log P=\log R\cos\tfrac12\varphi$\hfill .\,.\,.} & $8{,}782\,4527$\\[.3em]
\makebox[17em][l]{$\log Q=\log R\sin\tfrac12\varphi$\hfill .\,.\,.} & $7{,}877\,8355$\\[.3em]
\hline
$\text{Hinc }\tfrac12\varphi=$ & $7^0\,6'\,0''{,}935$\\[.3em]
$\varphi=$ & $14\quad12\quad1{,}87$\\[.3em]
\makebox[17em][l]{$\log R$\hfill .\,.\,.\,.\,.\,.\,.\,.\,.\,.\,.\,.\,.\,.\,.} & $8{,}785\,7960$\\[.3em]
\end{tabular}
\end{center}

\begin{center}\small Ad calculum confirmandum\end{center}

\begin{center}
\begin{tabular}{@{}ll@{}}
\makebox[17em][l]{$\tfrac12\log2\cos f$\hfill .\,.\,.\,.\,.\,.\,.\,.\,.\,.\,.} & $0{,}150\,0394$\\[.3em]
\makebox[17em][l]{$\tfrac12\log(l+x)=\log\dfrac m y$\hfill .\,.\,.\,.} & $8{,}635\,7566$\\[.3em]
\hline
 & $8{,}785\,7960$\\[.3em]
\end{tabular}
\end{center}

\SourceLeaf{134}

\begin{center}
\begin{tabular}{@{}ll@{}}
\makebox[17em][l]{$\tfrac12\log rr'$\hfill .\,.\,.\,.\,.\,.\,.\,.\,.\,.\,.\,.\,.\,.} & $0{,}326\,4939$\\[.3em]
\makebox[17em][l]{$\log\sin f$\hfill .\,.\,.\,.\,.\,.\,.\,.\,.\,.\,.\,.\,.\,.\,.} & $8{,}820\,2909$\\[.3em]
\makebox[17em][l]{$\mathrm{C.}\log\sin g$\hfill .\,.\,.\,.\,.\,.\,.\,.\,.\,.\,.\,.} & $1{,}262\,1765$\\[.3em]
\hline
\makebox[17em][l]{$\log b$\hfill .\,.\,.\,.\,.\,.\,.\,.\,.\,.\,.\,.\,.\,.\,.\,.\,.} & $0{,}408\,9613$\\[.3em]
\makebox[17em][l]{$\log\cos\varphi$\hfill .\,.\,.\,.\,.\,.\,.\,.\,.\,.\,.\,.\,.\,.} & $9{,}986\,5224$\\[.3em]
\hline
\makebox[17em][l]{$\log p$\hfill .\,.\,.\,.\,.\,.\,.\,.\,.\,.\,.\,.\,.\,.\,.\,.\,.} & $0{,}395\,4837$\\[.3em]
\makebox[17em][l]{$\log a$\hfill .\,.\,.\,.\,.\,.\,.\,.\,.\,.\,.\,.\,.\,.\,.\,.\,.} & $0{,}422\,4389$\\[.3em]
\end{tabular}
\end{center}

\begin{center}
\begin{tabular}{@{}ll@{}}
\makebox[17em][l]{$\log k$\hfill .\,.\,.\,.\,.\,.\,.\,.\,.\,.\,.\,.\,.\,.\,.\,.\,.} & $3{,}550\,0066$\\[.3em]
\makebox[17em][l]{$\tfrac32\log a$\hfill .\,.\,.\,.\,.\,.\,.\,.\,.\,.\,.\,.\,.\,.\,.\,.} & $0{,}633\,6584$\\[.3em]
\hline
 & $2{,}916\,3482$\\[.3em]
\makebox[17em][l]{$\log t$\hfill .\,.\,.\,.\,.\,.\,.\,.\,.\,.\,.\,.\,.\,.\,.\,.\,.} & $1{,}341\,1160$\\[.3em]
\hline
 & $4{,}257\,4642$\\[.3em]
\end{tabular}
\end{center}

\noindent Est itaque motus medius diurnus $=824''{,}7990$. Motus medius intra tempus $t=18091''{,}07=5^0\,1'\,31''{,}07$.

\begin{center}
\begin{tabular}{@{}ll@{}}
\makebox[17em][l]{$\log\sin\varphi$\hfill .\,.\,.\,.\,.\,.\,.\,.\,.\,.\,.\,.\,.\,.} & $9{,}389\,7262$\\[.3em]
\makebox[17em][l]{$\log206265$\hfill .\,.\,.\,.\,.\,.\,.\,.\,.\,.\,.\,.} & $5{,}314\,4251$\\[.3em]
\hline
\makebox[17em][l]{$\log e\text{ in secundis}$\hfill .\,.\,.\,.\,.\,.\,.\,.} & $4{,}704\,1513$\\[.3em]
\makebox[17em][l]{$\mathop{\mathrm{lo}\bar{\mathrm{g}}}\nolimits\sin E$\hfill .\,.\,.\,.\,.\,.\,.\,.\,.\,.\,.\,.\,.\,.} & $9{,}800\,0767_n$\\[.3em]
\makebox[17em][l]{$\log\sin E'$\hfill .\,.\,.\,.\,.\,.\,.\,.\,.\,.\,.\,.\,.} & $9{,}734\,4714_n$\\[.3em]
\hline
\makebox[17em][l]{$\log e\sin E$\hfill .\,.\,.\,.\,.\,.\,.\,.\,.\,.\,.\,.\,.} & $4{,}504\,2280_n$\\[.3em]
\makebox[17em][l]{$\log e\sin E'$\hfill .\,.\,.\,.\,.\,.\,.\,.\,.\,.\,.\,.} & $4{,}438\,6227_n$\\[.3em]
\end{tabular}
\end{center}

\[
\begin{aligned}
e\sin E&=-31932''{,}14=-8^0\,52'\,12''{,}14\\
e\sin E'&=-27455{,}08=-7\quad37\quad35{,}08
\end{aligned}
\]
\begin{center}Hinc anomalia media\end{center}

\begin{center}
\begin{tabular}{@{}ll@{}}
$\text{pro loco primo}=$ & $329^0\,44'\,27''{,}67$\\[.3em]
$\text{pro secundo}=$ & $334\quad45\quad58{,}73$\\[.3em]
\hline
$\text{Differentia}=$ & $5\quad1\quad31{,}06$\\[.3em]
\end{tabular}
\end{center}

II. In exemplo altero fit $f=31^0\,27'\,38''{,}32$, $\omega=-21'\,50''{,}565$, $l=0{,}086\,3565\,9$, $\log mm=9{,}353\,0651$, $\dfrac{mm}{\frac56+l}$ sive valor approximatus ipsius $h=0{,}245\,1454$; huic in tabula II respondet $\log yy=0{,}172\,2663$, unde deducitur $\dfrac{mm}{yy}=0{,}151\,6347\,7$, $x=0{,}065\,2781\,8$, hinc e tabula III sumitur $\xi=0{,}000\,2531$. Quo valore adhibito prodeunt valores correcti $h=0{,}245\,0779$, $\log yy=0{,}172\,2303$, $\dfrac{mm}{yy}=0{,}151\,6473\,5$, $x=0{,}065\,2907\,6$, $\xi=0{,}000\,2532$. Quodsi cum hoc valore ipsius $\xi$, unica tantum unitate in figura septima a priori div\={e}rso, calculus denuo repeteretur, $h$, $\log yy$, $x$ mutationem sensibilem non acciperent, quamobrem valor inventus ipsius $x$ iam est verus, statimque inde ad determinationem elementorum progredi licet. Cui hic non immoramur, quum nihil ab exemplo praecedente differat.

III. Haud abs re erit, etiam casum alterum ubi $\cos f$ negativus est exemplo illustrare. Sit $v'-v=224^0\,0'\,0''$, sive $f=112^0\,0'\,0''$, $\log r=0{,}139\,4892$, $\log r'=0{,}397\,8794$, $t=206{,}80919$ dies. Hic invenitur $\omega=+4^0\,14'\,43''{,}82$, $L=1{,}894\,2298$, $\log MM=0{,}672\,4333$, valor primus approximatus ipsius $\log H=0{,}646\,7603$, unde per solutionem aequationis $15^*$ obtinetur $Y=1{,}591\,432$, ac dein $x=0{,}037\,037$, cui respondet, in tabula III, $\xi=0{,}000\,0801$. \SourceLeaf{135}Hinc oriuntur valores correcti $\log H=0{,}646\,7931$, $Y=1{,}591\,5107$, $x=0{,}037\,2195$, $\xi=0{,}000\,0809$. Calculo cum hoc valore ipsius $\xi$ denuo repetito prodit $x=0{,}037\,2213$, qui valor, quum $\xi$ inde haud mutata prodeat, nulla amplius correctione indiget. Invenitur dein $\tfrac12 g=11^0\,7'\,25''{,}40$, atque hinc perinde ut in exemplo I

\begin{center}
\begin{tabular}{@{}ll@{}}
$\tfrac12(F-G)=$ & $3^0\,33'\,53''{,}62$\\[.3em]
$\tfrac12(F+G)=$ & $8\quad26\quad6{,}38$\\[.3em]
$F=$ & $12\quad0\quad0{,}00$\\[.3em]
$v=$ & $-100\quad0\quad0{,}00$\\[.3em]
$v'=$ & $+124\quad0\quad0{,}00$\\[.3em]
$G=$ & $4\quad52\quad12{,}76$\\[.3em]
$E=$ & $-17\quad22\quad38{,}04$\\[.3em]
$E'=$ & $+27\quad7\quad3{,}56$\\[.3em]
\end{tabular}
\end{center}

\begin{center}
\begin{tabular}{@{}ll@{}}
\makebox[17em][l]{$\log P=\log R\cos\tfrac12\varphi$\hfill .\,.\,.\,.} & $9{,}970\,0508$\\[.3em]
\makebox[17em][l]{$\log Q=\log R\sin\tfrac12\varphi$\hfill .\,.\,.\,.} & $9{,}858\,0553$\\[.3em]
$\tfrac12\varphi=$ & $37^0\,41'\,34''{,}27$\\[.3em]
$\varphi=$ & $75\quad23\quad8{,}54$\\[.3em]
\makebox[17em][l]{$\log R$\hfill .\,.\,.\,.\,.\,.\,.\,.\,.\,.\,.\,.\,.\,.\,.\,.\,.} & $0{,}071\,7097$\\[.3em]
\end{tabular}
\end{center}

\begin{center}\small Ad calculi confirmationem eruitur\end{center}

\begin{center}
\begin{tabular}{@{}ll@{}}
\makebox[17em][l]{$\log\dfrac M Y\BareRoot{(-2\cos f)}$\hfill .\,.\,.\,.\,.\,.\,.} & $0{,}071\,7097$\\[.3em]
\end{tabular}
\end{center}

\noindent In orbitis tam excentricis angulus $\varphi$ paullulo exactius computatur per formulam $19^*$, quae in exemplo nostro dat $\varphi=75^0\,23'\,8''{,}57$; excentricitas quoque $e$ maiori praecisione determinatur per formulam $1-2\sin(45^0-\tfrac12\varphi)^2$ quam per $\sin\varphi$; secundum illam fit $e=0{,}967\,6463\,0$.

Per formulam 1 porro invenitur $\log b=0{,}657\,6611$, unde $\log p=0{,}0595967$, $\log a=1{,}255\,7255$, atque logarithmus distantiae in perihelio
\[
\begin{aligned}
{}={}&\log\dfrac p{1+e}=\log(a(1-e))\\
{}={}&\log(b\operatorname{tang}(45^0-\tfrac12\varphi))=9{,}765\,6496.
\end{aligned}
\]

In orbitis tantopere ad parabolae similitudinem vergentibus loco epochae anomaliae mediae assignari solet tempus transitus per perihelium; intervalla inter hoc tempus atque tempora duobus locis propositis respondentia determinari poterunt ex elementis cognitis per methodum in art. 41 traditam, quorum differentia vel summa (prout perihelium vel extra duo loca proposita iacet vel intra), quum consentire debeat cum tempore $t$, calculo confirmando inserviet. — Ceterum numeri huius tertii exempli superstructi erant elementis in exemplo artt. 38 et 43 suppositis, quin adeo istud ipsum exemplum locum nostrum primum suppeditaverat: differentiae leviusculae elementorum hic erutorum unice a limitata praecisione tabularum logarithmicarum et trigonometricarum originem traxerunt.
